\abstractheading
  {A method to determine the minimal null control time of one-dimensional linear hyperbolic systems}
  {Guillaume Olive}
  {\affilJagiellonian}
  {Contributed talk}

How fast can a one-dimensional transport process be brought to rest by acting on
just one end? This is a question about the \emph{minimal null control time}, a
central notion in control theory that, for coupled hyperbolic systems, turns out
to be delicate.

Here, we consider first-order linear hyperbolic systems of the form:
\[
\left\{
\begin{aligned}
\frac{\partial y}{\partial t}(t,x)
  +\Lambda(x)\frac{\partial y}{\partial x}(t,x)
  &=M(x)y(t,x),\\
y_{-}(t,1)&=u(t),
  &y_{+}(t,0)&=Qy_{-}(t,0),\\
y(0,x)&=y^{0}(x),
\end{aligned}
\right.
\qquad (t,x)\in(0,T)\times(0,1).
\]

The components propagate at the speeds prescribed by $\Lambda$, interact
through the internal coupling $M$, are coupled at the boundary through $Q$,
and are steered by a control $u$ acting on one side. Such systems arise as
linearizations, around steady states, of the nonlinear balance laws governing
many transport phenomena [BC16; YK22]. Already in this linear setting,
determining the minimal control time is a subtle problem, as it depends not only
on the propagation speeds but on the couplings.

In this talk, we present a general method that computes this minimal control time
exactly, whenever $\Lambda$ and $M$ are sufficiently smooth functions of
$x$. This presentation is based on the joint work [HO25] with Long Hu.

\begin{abstractreferences}
  \referenceitem{BC16}{G. Bastin and J.-M. Coron. \emph{Stability and boundary
    stabilization of 1-D hyperbolic systems}. Vol.~88. Progress in Nonlinear
    Differential Equations and their Applications. Subseries in Control.
    Birkhäuser/Springer, [Cham], 2016, pp.~xiv+307.}
  \referenceitem{HO25}{L. Hu and G. Olive. ``A method to determine the minimal
    null control time of 1D linear hyperbolic balance laws''. In:
    \emph{J. Differential Equations} 440. part~1 (2025), pp.~113--455.}
  \referenceitem{YK22}{H. Yu and M. Krstic. \emph{Traffic Congestion Control by
    PDE Backstepping}. Systems \& Control: Foundations \& Applications.
    Birkhäuser Cham, 2022, pp.~XVII, 356.}
\end{abstractreferences}
